\documentclass{article}
\usepackage[utf8]{inputenc}
\usepackage[T1]{fontenc}
\usepackage{geometry}
\usepackage{graphicx}
\usepackage{float}
\usepackage{hyperref}
\usepackage[english,indonesian]{babel}

\title{Optimasi Pemrosesan Template LaTeX Menggunakan Pendekatan Berbasis AI}
\author{Academic Support Agent}
\date{\today}

\begin{document}
\maketitle

\begin{abstract}

\end{abstract}

\chapter{PENDAHULUAN}


\chapter{TINJAUAN PUSTAKA}


% Original content (commented out):
% \documentclass[11pt,letterpaper]{article}
\usepackage[lmargin=1in,rmargin=1in,tmargin=1in,bmargin=1in]{geometry}
\usepackage{hw}

% -------------------
% Content
% -------------------
\begin{document}
\homework{1}{MM/DD}

% Question
\prob This is an unnumbered problem. 
	\ansbox{2.5cm}


% Question 1
\problem This is a numbered problem.
	\ansbox{2.5cm}


% Question 2
\problem This is the second numbered problem.
	\ansbox{2.5cm}


% Question 3
\pointproblem{5} This problem has several parts:
	\begin{enumerate}[(a)]
	\item The first part.
	\item The second part. 
	\item The third part. 
	\end{enumerate}
\ansbox{6cm}















\end{document}

\end{document}